\section{Choosing a Winner: The Two Comparison Metrics}

Once every candidate design has been checked against both safety gates,
many designs pass. The question becomes: among the safe ones, which is
actually best? ``Best'' needs a definition, and this project deliberately
tried two different ones rather than assuming a single answer.

\subsection{Metric 1: Maximum Packing Density}

Simple: among every design that clears both safety gates, pick the one
with the most holes per unit area. More holes per area means more open
area overall (for a fixed hole size), which means more light gets through
--- directly useful once transmission/brightness is brought back into the
picture later.

\intuitbox{Why not just rank by \emph{safety margin} instead? Because a design that is
safe with enormous margin (holes spaced very far apart) is not necessarily
useful --- it might only fit a handful of holes. Ranking by margin alone
would reward uselessly sparse designs. Density is the metric that actually
reflects the practical goal.
}
\subsection{Metric 2: The Pareto Front}

Metric 1 collapses everything to a single winning point, which hides any
trade-off that might exist between density and safety margin. Metric 2
instead reports the full \emph{Pareto front} --- every design for which no
other design is simultaneously denser \emph{and} safer. This shows the
whole trade-off curve rather than one point on it.

In the current dataset, this front collapses to a single point at every
myopia level --- meaning, at the tested spacing resolution, safety in this
project turned out to be close to binary (trivially safe versus clearly
unsafe) rather than a graded trade-off. This was checked to be a genuine
property of the data, not a bug in the Pareto logic itself, using a
synthetic test dataset built specifically to confirm the dominance-checking
code \emph{can} produce a multi-point front when one actually exists.

\codemapbox{Both metrics live in \texttt{comparison\_metric/metrics.py}, operating on
whatever survives \texttt{gates.py}'s filter --- deliberately kept as two
separate, swappable functions so a future third metric does not require
rewriting the comparison logic.
}
\section{Closing the Loop: Confirming the Theory Against Real Physics}

Everything above --- the ghost formula, the vignetting geometry, both
metrics --- uses fast, \emph{closed-form} formulas. They were validated once
against real wave-optics simulation for a simple two-hole case
(Section~2), but never checked against the specific, sometimes irregular,
multi-hole patterns that actually won the comparison. A formula validated
for a simple case is not automatically guaranteed to hold for a more
complex one built from the same physics.

\subsection{The Confirmation Step}

The winning designs from Metric~1 were each propagated through the
project's full multi-hole wave-optics engine (the same validated
simulation machinery used for the original single-aperture optimization,
not a shortcut) --- and the resulting simulated image was inspected
directly, with actual peak-detection on the simulated retina, rather than
trusting the formula's prediction.

\intuitbox{This distinction matters: the closed-form formula \emph{predicts} where
ghost images should land. The full simulation \emph{measures} where they
actually land, by literally propagating light through the eye model and
looking at the resulting image. Confirmation means these two independent
methods agree; it is not simply re-running the same formula a second time.
}
The result: \textbf{9 out of 9 finalist designs were confirmed} --- the
full simulation reproduced each finalist's predicted ghost-safety and
vignetting-safety exactly, at every tested gaze angle. One genuinely useful
side-finding also came out of this step: three of the nine finalists
revealed that an older reference file (\texttt{blur\_sweep.csv}) had been
computed at a since-corrected, lower simulation resolution --- caught
because three finalists with \emph{different} hole sizes reported the
\emph{identical} stale blur value to many decimal places, which is not
something real, correctly-computed diffraction physics would ever do for
different apertures. That reference file has since been regenerated at
the corrected resolution.

\codemapbox{This confirmation pass lives in
\texttt{array\_geometry\_lab\_phase2\_confirmation/}, importing (not
copying) the original validated simulation engine directly --- correctness
matters more than speed here, since only a handful of finalist designs
need checking, unlike the thousand-plus-candidate sweep that justified
using fast closed-form copies everywhere else.
}
\section{Summary: What a Design Needs to Pass}

\begin{center}
\begin{tabular}{cl}
\toprule
\textbf{Step} & \textbf{Requirement} \\
\midrule
1 & Hole size fixed by the existing acuity-optimal chart, per myopia level \\
2 & Generate a candidate geometry (pattern, hole count, spacing) \\
3 & \textbf{Gate 1:} no ghost-image pair violates the resolvable-separation limit \\
4 & \textbf{Gate 2:} Gate 1 holds at every tested gaze angle, not just $\varphi=0$ \\
5 & Among survivors, rank by packing density (Metric 1) or inspect the trade-off front (Metric 2) \\
6 & Confirm the winner(s) against full wave-optics simulation, not just the formula \\
\bottomrule
\end{tabular}
\end{center}

A design that passes all six steps has been checked, independently, at two
different levels of physical detail (fast formula and full simulation), and
is the point at which this project's theoretical work hands off to physical
fabrication and measurement.

\end{document}
